% a06_pgate_a: Step 3(F), fifth part (input from sec/a06_BALtorus_ct8.tex): the new fluctuation field a of Ba2 (2.88)-(2.89)
% at step k (value, derivative and Hoelder bounds) and the geometry of the g_k-gate. Continued in a06_pgate_b (the gated terms
% of P^(k)) and a06_pgate_c (the final step, Ba2 (2.116)).

\subsubsection*{The $g_k$-gated terms of $P^{(k)}$ and the final step \cite[p.~579, (2.116)]{Ba2}}

\emph{Setting.} Step $k$, $1\le k\le K-1$, in scale-$k$ units, with the Setting of \cref{thm:a06f:C2} ($\Omega_A$, $\Omega_1$, $\Omega'_1=B^k(\Lambda^{(k)}_7)$, $\mathrm{Col}=\mathrm{Col}_k$, $\Phi=\phi^{(k)}$, $p=p(t_k)$, $r=r(t_k)$). At step $k+1$ \cite{Ba2} writes $\theta_{k+1}A^{(k)}=\theta_{k+1}A'^{(k)}+\tilde B^{(k)}+\theta_{k+1}B^{(k+1),\eta}$ \cite[(2.88)--(2.89)]{Ba2}, with $A'^{(k)}=a_k\zeta^{(k)}G_kQ_k^*A'$ (\cref{lem:a06d:C62}), $\theta_{k+1}$ the smooth cutoff of $B^k(\Lambda^{(k)}_2)$ changing in a slice of thickness $<M$ \cite[p.~567]{Ba2}, and every term of $\tilde B^{(k)}$ carrying the factor $\theta_{k+1}$. Put
\[
\tilde B\defeq(1-\theta_{k+1})\theta_kA^{(k)}+\theta_{k+1}B^{(k+1),\eta},\qquad a\defeq\theta_{k+1}A'^{(k)}+\tilde B^{(k)},\qquad\text{so that}\qquad\theta_kA^{(k)}=\tilde B+a,
\]
since $\theta_k=1$ on $\supp\theta_{k+1}\subset B^k(\Lambda^{(k-1)'}_7)\subset B^{k-1}(\Lambda^{(k-1)}_2)$ (\cref{lem:a06d:G3}); $\tilde B$ is the background of \cite[p.~582]{Ba2}, regular by \cite[(2.94)--(2.98)]{Ba2}, and $\supp a\subset\supp\theta_{k+1}$, which lies within $M$ of $B^k(\Lambda^{(k)}_2)$, hence in $B^k(\Lambda^{(k)}_1)$ ($\dist(\Lambda^{(k)}_2,\Lambda^{(k)c}_1)\ge r-1>M+1$ under the caps of Step~3(D$_2$)). The gate is $g_k$ ($\supp g_k\subset\Lambda^{(k)}_7$, $g_k=1$ on $\Lambda^{(k)}_8$, $|\partial^\eta g_k|\le c_g$, as in \cref{prop:a06f:H}). The cutoffs are taken, as the printed ``smooth'' allows, from fixed classes: $\theta_{k+1}$ a fixed smooth function of the mollified distance to $B^k(\Lambda^{(k)}_2)$, so that $|\partial^\eta\theta_{k+1}|\le c_\theta$ and $|\partial^\eta\partial^\eta\theta_{k+1}|\le c_{\theta2}$ uniformly in $k$ and $\eta$ (a smooth function of the lattice sup-distance would have second differences of order $\eta^{-1}$ along the edges of the slice); and $\zeta^{(k)}$ with $|\zeta^{(k)}|\le1$, $|\partial^\eta_x\zeta^{(k)}|\le1$, $|\partial^\eta_x\partial^\eta_x\zeta^{(k)}|\le2$ (the cut-off of \cref{thm:a06:BALA}(b) is so chosen: a fixed smooth profile of the mollified distance to $y$, as for $\theta_{k+1}$, over its transition region of width $\frac15r-2M$, which is $\ge2$ for $r\ge10M+10$, a cap on $\eps_0$ free of $k$). Put $\delta'\defeq\min(\delta_*,\delta_0)$, $S'_*\defeq\sum_{n\in\Z^3}e^{-\delta'|n|/2}$, and $C'(\delta')$ for the constant $C'$ of \cref{lem:a06d:C62} with $\delta'$ in place of $\delta_*$.

\begin{lemma}[the field $a$]\label{lem:a06g:A}
On $\supp\theta_{k+1}$, for $t_k\le t_0$, where $t_0\in(0,1]$ is fixed by $C_\kappa t_0^\kappa\le C_{2.91}b_0$, $C'_\kappa t_0^\kappa\le b_0$ and $C^H_\kappa t_0^\kappa\le b_0$ ($C_\kappa$, $C'_\kappa$ the constants of (R2) of Step~3(C$_2$), $C^H_\kappa$ that of (c) below; a cap on $\eps_0$, free of $k$; $p(t)\ge b_0$): \textup{(a)} $|a|\le C_{2.91}p+C_\kappa t_k^\kappa\le2C_{2.91}p$; \textup{(b)} for every fine bond $b=(x,\mu)$ with $x$ or $x+\eta e_\mu$ in $\supp\theta_{k+1}$, $|\partial^\eta_\mu a(x)|\le C'_{2.91}p+C'_\kappa t_k^\kappa\le(C'_{2.91}+1)p$, with
\[
C'_{2.91}\defeq c_\theta C_{2.91}+a(C_1e^{\delta_*}+c_0e^{\delta_0})S'_*(C_{2.86}+C'(\delta'));
\]
\textup{(c)} for parallel bonds $b,b'$ within distance $M$ of $\supp\theta_{k+1}$ with $|b-b'|\le1$, $|\partial^\eta a(b)-\partial^\eta a(b')|\le C^H_{2.91}p\,|b-b'|^{\alpha_0}$. The constants are free of $k$, $\eps$, $\eps_0$, the torus and the history. So $a$ satisfies the hypotheses on an external vector field of \cite[p.~412]{Ba3} with $\norm{\cdot}_{1,\alpha_0}$-norm $\le C_ap$, $C_a\defeq2C_{2.91}+C'_{2.91}+1+C^H_{2.91}$.
\end{lemma}

\begin{proof}
(a) is \cref{lem:a06d:C62} for $\theta_{k+1}A'^{(k)}$ and (R2) of Step~3(C$_2$) for $\tilde B^{(k)}$ \cite[(2.90)]{Ba2}; $C_\kappa t^\kappa\le C_{2.91}p$ for $t\le t_0$.

(b) \emph{Lattice distances.} Fine points are at mutual distances in $\eta\Z$, and the slice has thickness $<M$, so $\supp\theta_{k+1}$ lies within $M-\eta$ of $B^k(\Lambda^{(k)}_2)$ and both $x$ and $x+\eta e_\mu$ lie within $M$ of it; \cref{lem:a06d:C62} and its proof apply at both points (the proof bounds $A'^{(k)}$ itself there, without the factor $\theta_{k+1}$). \emph{Leibniz rule.} $\partial_\mu(\theta_{k+1}A'^{(k)})(x)=(\partial_\mu\theta_{k+1})(x)A'^{(k)}(x+\eta e_\mu)+\theta_{k+1}(x)(\partial_\mu A'^{(k)})(x)$. The first term is $\le c_\theta C_{2.91}p$. In the second,
\[
\partial_\mu A'^{(k)}(x)=a_k\sum_y\big[(\partial_{\mu,x}\zeta^{(k)})(x,y)(G_kQ_k^*)(x+\eta e_\mu,y)+\zeta^{(k)}(x,y)(\partial_\mu G_kQ_k^*)(x,y)\big]A'(y);
\]
both kernels are $\le(C_1e^{\delta_*}+c_0e^{\delta_0})e^{-\delta'|x-y|}$, by $|\partial\zeta|\le1$, \cref{thm:a06:A5} ($|\cdot|_1\ge|\cdot|$ and one bond shift) and the derivative clause of \cite[Prop.~2.2, (2.58)]{Ba2} (\cite[Thm., (1.9)--(1.10)]{Ba4}, at zero background on the torus, where the depth restriction is void, \cref{lem:a06A:torus}); both vanish unless $|x-y|<r-2M+\eta$, and such $y$ lie in $\Lambda^{(k)}_0$ by the layer count of \cref{lem:a06d:C62} ($2(r-c_b)>r-2M+M+1$). Split the $y$-sum as there: on $\Lambda^{(k)}_1$, \cref{lem:a06d:C61} gives $\le(\dots)S'_*C_{2.86}p$; on $\Lambda^{(k)}_0\setminus\Lambda^{(k)}_1$, $|x-y|\ge r-c_b-M-1$ and the factor $e^{-\delta'|x-y|}$ with $E(1,\delta'/2)$ turn $p/(\lmu t)$ into $C'(\delta')p$. With $|\theta_{k+1}|\le1$ and $a_k\le a$ the second term is $\le a(C_1e^{\delta_*}+c_0e^{\delta_0})S'_*(C_{2.86}+C'(\delta'))p$. For $\tilde B^{(k)}$ the derivative bound is (R2).

(c) Write $\partial(\theta_{k+1}A'^{(k)})=(\partial\theta_{k+1})\,\mathsf SA'^{(k)}+\theta_{k+1}\partial A'^{(k)}$ ($\mathsf S$ the bond shift) and telescope the difference at $b,b'$ term by term. The factors $\partial\theta_{k+1}$, $\theta_{k+1}$ and the kernel factor $\partial\zeta^{(k)}$ are Lipschitz on the $\eta$-lattice with constants $c_{\theta2}$, $c_\theta$ and $2$; at separation $|b-b'|\ge\eta$ a Lipschitz bound is $\le$ the constant times $|b-b'|^{\alpha_0}$, since $|b-b'|\le1$; their cofactors ($A'^{(k)}$, $\partial A'^{(k)}$ and the kernel sums) are bounded by \cref{lem:a06d:C62} and by (b), the kernel $(G_kQ_k^*)(\cdot+\eta e_\mu,y)$ of the $\partial\zeta^{(k)}$ term being itself Lipschitz in $x$ with constant $c_0e^{\delta_0}e^{-\delta_0\dist(\cdot,y)}$ by the derivative clause of \cite[(2.58)]{Ba2} (a difference along a lattice path of length $|b-b'|\le1$), summed over $y$ as in (b). The $\theta$-factors are telescoped first, so each cofactor is either evaluated at a bond of $\supp\theta_{k+1}$ or is a difference of two kernel sums of the type of (b) at points within $M+1$ of $B^k(\Lambda^{(k)}_2)$; the bounds of (b) hold there, since the kernel bounds are global on the torus and the layer count with $M+1$ in place of $M$ keeps the $y$-range in $\Lambda^{(k)}_0$. The term $\zeta^{(k)}(\partial G_kQ_k^*)A'$: its variation in $b$ is the H\"older clause of \cite[Prop.~2.1, (2.24)]{Ba1} (\cite[Thm., (1.10)]{Ba4}) for $\partial G_kQ_k^*\delta_y$, $|(\partial G_kQ_k^*)(b,y)-(\partial G_kQ_k^*)(b',y)|\le c_0|b-b'|^{\alpha_0}e^{-\delta_0\min(\dist(b,y),\dist(b',y))}$; sum over $y$ as in (b). For $\tilde B^{(k)}$: by \cite[(2.89)]{Ba2} each of its three terms is $\theta_{k+1}$ times a kernel sum of the same type ($\zeta^{(k)}G_kQ_k^*$, its tail $(1-\zeta^{(k)})G_kQ_k^*$, or $\zeta^{(k)}G_kQ_k^*$ composed with the covariance difference of \cref{lem:a06A:CLdiff}) applied to a block field bounded by $Cp/(\lmu t)$ on $\Lambda^{(k)'}_0$ (the restrictions \cite[(2.17), (2.55)]{Ba2}); the same computation bounds its H\"older quotient, and the separation $\ge r-2c_b-M$ between the evaluation points in $\supp\theta_{k+1}$ and the support of the field, through $E(1,\delta/2)$, gives $C^H_\kappa t_k^\kappa$ in place of $p$, as in (R2).
\end{proof}

\begin{lemma}[the gate]\label{lem:a06g:BAD}
Let $\mathrm{Bad}_k$ be the set of scale-$k$ unit cubes meeting $\supp((1-g_k)a)$. Then $\mathrm{Bad}_k\subset N_M(B^k(\Lambda^{(k)}_2))\setminus B^k(\Lambda^{(k)}_8)\subset\Omega_1$, every cube of $\mathrm{Bad}_k$ lies within $M$ of a cube of $\Lambda^{(k)}_2\cap\Lambda^{(k)c}_8$, and
\[
|\mathrm{Bad}_k|\le(2M+1)^3|\Lambda^{(k)}_2\cap\Lambda^{(k)c}_8|\le(2M+1)^3\big(|\Lambda^{(k)}_2\cap\Lambda^{(k)c}_7|+|\Lambda^{(k)}_7\setminus\Lambda^{(k)}_8|\big),\qquad\Lambda^{(k)}_2\cap\Lambda^{(k)c}_7\subset\mathrm{Col}_k.
\]
\end{lemma}

\begin{proof}
$(1-g_k)a$ vanishes outside $\supp\theta_{k+1}\subset N_M(B^k(\Lambda^{(k)}_2))$ and on $B^k(\Lambda^{(k)}_8)$, where $g_k=1$. Let $y\in\mathrm{Bad}_k$ and let $x$ be a cube of $B^k(\Lambda^{(k)}_2)$ within $M$ of $y$. If $y\in\Lambda^{(k)}_2$, then $y\in\Lambda^{(k)}_2\cap\Lambda^{(k)c}_8$ itself. If $y\notin\Lambda^{(k)}_2$, then $\dist(x,\Lambda^{(k)c}_2)\le M+1<6(r-1)\le\dist(\Lambda^{(k)}_8,\Lambda^{(k)c}_2)$, so $x\notin\Lambda^{(k)}_8$ and $x\in\Lambda^{(k)}_2\cap\Lambda^{(k)c}_8$. A cube of $\Lambda^{(k)}_2\cap\Lambda^{(k)c}_8$ has at most $(2M+1)^3$ cubes within $M$ of it. The inclusions: $\Lambda^{(k)}_8\subset\Lambda^{(k)}_7$, and $\Lambda^{(k)}_2\subset\Lambda^{(k)}_0\subset\Lambda^{(k-1)'}_7$ \cite[pp.~566, 570]{Ba2}, so $\Lambda^{(k)}_2\cap\Lambda^{(k)c}_7\subset\Lambda^{(k-1)'}_7\cap\Lambda^{(k)c}_7=\mathrm{Col}_k$ and $N_M(B^k(\Lambda^{(k)}_2))\subset B^k(\Lambda^{(k)}_1)\subset\Omega_1$.
\end{proof}
